% Werke X/1 (1917), native171 lower–173, printed164 lower–166.
\begingroup
\setlength{\headheight}{25pt}
\ArithmeticHead{164}{L}{ANALYSIS.\quad NACHLASS.}
\begin{center}\small BEMERKUNGEN ZU DEN ABSCHNITTEN [II.] UND [III.].\end{center}
\par\indent Das in [II.] [1.] auftretende Integral und seine Umkehrung hatte \textsc{Gauss}, wie die Tagebuchnotiz Nr. 32 vom 9. September 1796 und art. [5.] der \textit{Exercitat. mathem.} zeigen (siehe oben S. 140 dieses Bandes), anfangs neben dem lemniskatischen Integrale betrachtet. Auffallend ist dort wie auch hier die Benutzung des Buchstaben \(x\) in zweierlei Bedeutung, besonders in der letzten Formel des Artikels, wo überdies \({}^3\!\surd(1-x^3)\) ein Schreibfehler für \(\surd(1-x^3)\) ist. Die Reihen, die zwischen den beiden Quotienten von unendlichen Produkten stehen, deuten auf eine Untersuchung der Konvergenz der Produkte; vergl. auch [II.] [5.].

In [II.] [2.] und [3.] erscheint die Funktion \(\sin\text{ lemn }x\) als Quotient, das einemal von doppelt unendlichen Produkten, das anderemal von beständig konvergenten Potenzreihen. Das hier sowie in [II.] [5.] und [6.] auftretende \(\Pi\) hat die Bedeutung
\[\Pi=\frac{\varpi}{2}=\int_0^1\frac{dp}{\sqrt{1-p^4}}.\]

In [II.] [5.] werden die Nullstellen der Funktion sinus lemniscaticus schematisch dargestellt; die Reihen \(\sum\frac1{\rho^4}\) und
\[\sum\frac{m^4-6m^2n^2+n^4}{(m^2+n^2)^4}=\tfrac12\sum\left\{\frac1{(m+n\surd-1)^4}+\frac1{(m-n\surd-1)^4}\right\}\]
deuten auch hier (vergl. [II.] [1.]) auf Konvergenzuntersuchungen\footnote{Ein Ausdruck, der an die letztere Reihe erinnert, wird in der Notiz Nr. 61 des \textit{Tagebuchs} (19.--21. März 1797) mit \(\Pi\) in Verbindung gebracht.}.
\clearpage
\SourceLeaf{172}\fancyhead{}\fancyhead[R]{\small165}\fancyhead[C]{\scriptsize BEMERKUNGEN ZU DEN ABSCHNITTEN II. UND III.}
Auch [II.] [4.] gehört noch in diesen Zusammenhang; wir finden hier die schon wiederholt betrachtete Potenzreihe für \(\sin\text{ lemn }x\); dann aber wird die Zerfällung des Wertes dieser Funktion für ein \textls[120]{komplexes Argument} mit Hilfe des Additionstheorems und der Formeln
\[\begin{aligned}
\sin\text{ lemn }u\surd-1&=\surd-1\sin\text{ lemn }u,\\
\cos\text{ lemn }u\surd-1&=\frac1{\cos\text{ lemn }u}
\end{aligned}\]
(vergl. Werke III, S. 411) gegeben. Übersichtlicher angeordnet als für die Funktion sinus lemniscaticus folgen in [II.] [6.] die analogen Untersuchungen für cosinus lemniscaticus.

Die in [II.] [3.] ausgeführten Zahlenrechnungen werden in [II.] [7.] mit größerer Schärfe wiederholt. Die in [II.] [7.] und [8.] von \textsc{Gauss} benutzten Zeichen \(M,N,\mu,\nu\) bedeuten die Zähler und Nenner von \(\sin\text{ lemn}\) und von \(\cos\text{ lemn}\), also
\[\begin{aligned}
M&=x-\frac1{16}x^5-\frac1{10080}x^9+\frac{23}{259459200}x^{13}+\ldots,\\
N&=1+\frac1{12}x^4-\frac1{10080}x^8+\frac{17}{19958400}x^{12}+\ldots,\\
\mu&=1-\frac12x^2-\frac1{24}x^4-\frac1{240}x^6-\ldots,\\
\nu&=1+\frac12x^2-\frac1{24}x^4+\frac1{240}x^6-\ldots;
\end{aligned}\]
dieselben Funktionen hat \textsc{Gauss} später (vergl. Werke III, S. 405, 406) mit \(P,Q,p,q\) bezeichnet; diese Bezeichnungen werden uns bei irregulären (vergl. S. 170, Fußnote) Leisteaufzeichnungen weiter unten im Abschnitt [IV.], art. [1.], [2.] noch begegnen.

Die Reihen für \(\mu\) und \(\nu\) hat \textsc{Gauss} mit der Bezeichnung \(p,q\) auch auf das vordere Schutzblatt des I. Bandes seines Exemplars von \textsc{Schulzes} \textit{Tafeln} eingetragen.

In [II.] [3.] und [7.] wird zunächst der Wert
\[a=M\frac{\varpi}{2}=N\frac{\varpi}{2}\]
mit Hülfe der Potenzreihen für \(M\) und \(N\) berechnet. Auf vier Dezimalstellen hatte \textsc{Gauss} diese Zahl (1,2453) auch auf die Einbanddecke des \textsc{Leiste} hingeschrieben (siehe oben [I.] [1.], S. 145). Es werden dann in [II.] [7.] und [8.] die Ausdrücke von \(N2\varphi\) und \(M2\varphi\) durch \(M\varphi\) und \(N\varphi\), ferner die Darstellung von \(N\lambda\varpi\) für ein beliebiges positives ganzzahliges \(\lambda\) als Potenz von \(a\) gegeben. Diese Resultate finden in der Notiz Nr. 63 des \textit{Tagebuchs} vom 29. März besondere Erwähnung. In [II.] [8.] werden dann die Werte \(M\frac{\varpi}{4}\) und \(N\frac{\varpi}{4}\) berechnet, und mit Hilfe der Formel für \(M2\varphi\) der in [II.] [3.] unmittelbar berechnete Wert \(M\frac{\varpi}{2}\) bestätigt. Die Formel für \(N2\varphi\) ergibt dann
\[N\varpi=2\left(M\frac{\varpi}{2}\right)^4,\]
und \textsc{Gauss} findet nun durch \textls[120]{numerische Induktion} die Gleichung
\[N\varpi=e^{\frac{\pi}{2}}.\]
Dieses Ergebnis wird in derselben Notiz Nr. 63 des \textit{Tagebuchs} als \textls[120]{maxime memorabile} bezeichnet.
\clearpage
\SourceLeaf{173}\ArithmeticHead{166}{L}{ANALYSIS.\quad NACHLASS.}
Die in [II.] [7.] und [8.] gegebenen Beziehungen zwischen den vier Funktionen \(M,N,\mu,\nu\) findet man auch vervollständigt auf S. 63 des Handbuchs 16, Bb (begonnen November 1801), abgedruckt Werke III, S. 410, 411; die Formeln von [II.] [10.] sind in dem Mai 1809 angefangenen Handbuch 19, Be, S. 80, 81 weiter ausgeführt, abgedruckt Werke III, S. 408.

In [III.] sind die auf die Lemniskatenteilung bezüglichen Leisteaufzeichnungen mit zwei andern Notizen\footnote{Diese sind zwar schon Werke III, S. 421, art. [13.] abgedruckt, der Abdruck ist jedoch unvollständig und durch unrichtige Deutung der von \textsc{Gauss} benutzten Zeichen entstellt.\ClosingRule} zusammengestellt, die aber der Handschrift nach zu urteilen alle aus derselben Zeit herrühren. Übrigens ist diese Zeit durch die Tagebuchnotizen Nr. 60 vom 19. März und Nr. 62 vom 21. März 1797 auf das genaueste festgelegt.

In [III.] [1.] sind in der Gleichung 48. Grades, die die Siebenteilung liefert, nur die Koeffizienten angedeutet. Die Gleichung lautet nach einer brieflichen Mitteilung von \textsc{K. Schwering}:
{\small\[16(1-x^4)\left(\frac{1-5x^4-5x^8+x^{12}}{1+20x^4-26x^8+20x^{12}+x^{16}}\right)^2=\left(\frac{3-6x^4-x^8}{1+6x^4-3x^8}\right)^2;\]}
\noindent sie ergibt sich, indem man in
\[\begin{array}{l@{\qquad\qquad}c}(\ast)&(\sin\text{ lemn }4\varphi)^2=(\sin\text{ lemn }3\varphi)^2\end{array}\]
beiderseits aus Abschnitt [I.], artt. [4.] und [6.] die Ausdrücke durch \(x=\sin\text{ lemn }\varphi\) einsetzt und nachher den Faktor \(x^2\) unterdrückt. Man erkennt dann, daß der Gleichung (\(\ast\)) die Bestimmung
\[7\varphi=2(m+m'\surd-1)\varpi\]
für ganzzahlige \(m,m'\) entspricht.

In [III.] [2.] bedeutet \(\sin\) wieder \(\sin\text{ lemn}\) und \(R\) steht (wie auch schon in [II.] [6.]) für \(90^0\), bedeutet also \(\frac{\varpi}{2}\). Daß es sich um die Fünfteilung der Lemniskate handelt, lehrt auch die kleine Figur.

In [III.] [3.] wird dann die \textls[120]{transzendente Bestimmung} der reellen Wurzeln der Teilungsgleichung vorgenommen. Die kleine Tabelle links gibt die Potenzen von \(\frac{\varpi}{2}\); eine kleinere solche Tabelle war schon oben [II.] [3.] vorgekommen, eine ähnliche, die aber mit viel größerer Schärfe gerechnet --- also jedenfalls später zusammengestellt ist ---, findet sich auf S. 3 der Scheda Aa (begonnen Juli 1798), abgedruckt Werke III, S. 413.

Die artt. [4.] und [5.] des Abschnitts [III.] sind durch die vom Herausgeber hinzugefügten Einschaltungen verständlich gemacht. Das »Determ. rad. imag.« (S. 162) findet seine Erledigung in der Tagebuchnotiz Nr. 60 vom 19. März 1797; vergl. auch Abschnitt [III.], art. [2.]. Der Schluß von [4.], von »\(S=\surd720-26\)« an, könnte etwas später eingetragen sein, vielleicht nach der in Nr. 63 des Tagebuchs aufgezeichneten Entdeckung, daß die Fünfteilung \textls[120]{geometrisch}, d. h. mit Lineal und Zirkel, ausführbar sei. Die \textls[120]{algebraische Darstellung} der Wurzeln findet sich nämlich im \textsc{Leiste} nicht, sondern erst auf den in [6.] und [7.] abgedruckten Zetteln.

Der Zusammenhang mit den Tagebuchnotizen Nr. 60---Nr. 63 zeigt, daß in den Abschnitten [II.] und [III.] zusammengestellten Aufzeichnungen Entdeckungen wiedergeben, die zwischen dem 19. und dem 29. März 1797 gemacht worden sind.
\begin{flushright}\textsc{Schlesinger.}\end{flushright}
\clearpage
\endgroup
